\documentclass[10pt,a4paper]{amsart}
\usepackage[margin=19mm]{geometry}
\usepackage{amsmath,amssymb,mathtools}
\usepackage[scr=rsfs,cal=euler]{mathalfa}
\usepackage{xfrac}
\usepackage[unicode,hidelinks,bookmarks=false]{hyperref}
\newcommand{\ie}{i.e.\ }
\newcommand{\cf}{cf.\ }
\newcommand{\resp}{respectively\ }
\newcommand{\Subsection}{\subsection}
\providecommand{\nobreakdash}{\nobreak\hbox{-}}
\setlength{\emergencystretch}{2em}
\multlinegap0pt
\usepackage{color}
\usepackage{amscd}
\usepackage{overpic}
\usepackage{tabularx}
\usepackage{bm}
\newtheorem{theorem}{Theorem}[section] 
\newtheorem{lemma}[theorem]{Lemma}
\newtheorem{corollary}[theorem]{Corollary}
\newtheorem{proposition}[theorem]{Proposition}
\newtheorem{claim}[theorem]{Claim}
\newtheorem{fact}[theorem]{Fact}
\newtheorem{definition}[theorem]{Definition}
\newtheorem{remark}[theorem]{Remark}
\newtheorem{remarks}[theorem]{Remarks}
\newtheorem{example}[theorem]{Example}
\newtheorem{notation}[theorem]{Notation}
\newtheorem{assertion}[theorem]{Assertion}
\newtheorem{problem}[theorem]{Problem}
\def\Bline{%
\noalign{\ifnum0=`}\fi\hrule \@height 1pt \futurelet
\reserved@a\@xhline}
\def\theequation{\thesection.\arabic{equation}}
\allowdisplaybreaks[3]
\renewcommand{\baselinestretch}{1.2}
\renewcommand{\labelenumii}{$(${\upshape\theenumii}$)$}
\renewcommand{\labelenumi}{$(${\upshape\theenumi}$)$}
\renewcommand{\figurename}{Figure.}
\newcommand{\textR}[1]{\textcolor{red}{#1}}
\newcommand{\textB}[1]{#1}
\newcommand{\textG}[1]{#1}
\newcommand{\hs}[1]{\hspace{#1pt}}
\newcommand{\R}{\mathbb{R}}
\newcommand{\N}{\mathbb{N}}
\newcommand{\Q}{\mathbb{Q}}
\newcommand{\Z}{\mathbb{Z}}
\newcommand{\C}{\mathbb{C}}
\newcommand{\pd}[2]{\dfrac{\partial#1}{\partial#2}}
\newcommand{\n}{\mathbf{n}}
\renewcommand{\v}{\mathbf{v}}
\newcommand{\x}{\mathbf{x}}
\newcommand{\p}{\mathbf{p}}
\newcommand{\e}{\mathbf{e}}
\renewcommand{\phi}{\varphi}
\renewcommand{\epsilon}{\varepsilon}
\newcommand{\wtilde}[1]{\widetilde{#1}}
\newcommand{\rank}{\mathrm{rank}}
\newcommand{\hess}{\mathrm{hess}}
\renewcommand{\geq}{\geqslant}
\renewcommand{\leq}{\leqslant}
\newcommand{\dsum}[2]{\displaystyle \sum_{#1}^{#2}}
\begin{document}
\title[Height and distance-squared functions on implicit plane curv]{Height and distance-squared functions on implicit plane curves}
\author{Masaru Hasegawa}
\date{}
\maketitle
\noindent\emph{Source and licence.} Height and distance-squared functions on implicit plane curves. Version arXiv:2608.19517v1 (CC BY 4.0). This material is distributed under the Creative Commons Attribution 4.0 International licence, http://creativecommons.org/licenses/by/4.0/, as recorded in the arXiv OAI licence metadata for this exact version and shown on the abstract page https://arxiv.org/abs/2608.19517v1. Full rights notice: licenses/arxiv-2608.19517-CC-BY-4.0.txt.\par\smallskip
\noindent\textit{Editorial scope.} Counted unit: the original \texttt{theorem} environment \texttt{thm:distance} at source lines 271--286 together with its complete proof at lines 288--336; the delivered document keeps source lines 164--337 so that every referenced label and every citation resolves inside it. Native editable source is the paper's own concatenated TeX; the AMS wrapper is editorial formatting only. No publisher class, upstream script or unrelated artwork is redistributed.
\begin{equation*}
%\label{eq:Lagrange}
F:\R^2\times\R\to\R,\quad F(\x,\lambda)=g(\x)+\lambda f(\x).
\end{equation*}
It is well known that $F$ is the Lagrange function and $\lambda$ is
called the Lagrange multiplier. 
The condition $\nabla F(\x,\lambda) = 0$ holds if and only if $\nabla g(\x)=-\lambda\nabla f(\x)$ and $f(\x)=0$. 
Therefore, if $\nabla f(\x)\ne 0$ and $\nabla g(\x) \ne 0$ then $\nabla F(\x,\lambda) = 0$ if and only if $f(\x)=0$ is tangent to $g(\x)=d$ at $\x$ for some $d$. 
The relation between contact geometry and singularity theory has been extensively studied (see \cite{IRRT2015} for example).
This suggests that more degenerate differential conditions on $F$ correspond to more degenerate contact between $f(\x)=0$ and $g(\x)=d$.  

In \cite{Hasegawa2017}, the author gave an answer to this question. 
For an implicit surface $M=\{\x\in\R^3|f(\x)=0\}$, where $f:\R^3\to\R$ is a smooth function, the author introduced two families of functions:
\begin{align*}
%\label{eq:impheight}
& H:\R^3\times\R\times S^2 \to \R,\quad H(\x,\lambda,\v)= \langle \x,\v \rangle+\lambda f(\x);\\
%\label{eq:impdistance}
& D:\R^3\times\R\times\R^3\setminus M \to\R,\quad D(\x,\lambda,\p) = |\x-\p|^2+\lambda f(\x). 
\end{align*}
The functions $H$ and $D$ measure the contact of $M$ with planes and spheres, respectively, and hence play the roles of families of height and distance-squared functions on $M$. 
Formulas for geometric invariants of implicit curves are well known (see \cite{Goldman2005} for example), while inflections and vertices of parametrized plane curves have been studied from the viewpoint of singularity theory (see \cite{DT2018,ST2017} for example).
It is natural to consider similar families of functions on implicit plane curves. 
We investigate these invariants from a different viewpoint:
the degeneracies of the associated Lagrange functions and the geometric
meaning of the corresponding Lagrange multipliers. 
We show that these degeneracies characterize geometric properties of implicit plane curves, such as inflections and vertices, through their contact with lines and circles.

In Section 2, we introduce families of height and distance-squared functions on implicit plane curves and characterize their singularities in terms of inflections and vertices. 
In Section 3, we apply the distance-squared family to the evolute of an
implicit plane curve and show, through the ordinary $3/2$-cusp, a
difference between the evolutes obtained from the implicit and
parametric representations of the curve.

The author would like to express his sincere gratitude to Toshizumi Fukui for fruitful suggestions. 

%%%%%%%%%%%%%%%%%%%%%%%%%%%%%%%%%%%%%%%%%%%%%%%%%%

\section{Families of height and distance-squared functions on implicit plane curves}

Let $f:\R^2\to\R$ be a smooth function, and let $C$ denote an implicit plane curve $\{(x,y)=\x\in\R^2|f(\x)=0\}$. The unit normal vector to $C$ is given by $\n=\nabla f/|\nabla f|$, where $\nabla f$ is the gradient as row vector. The curvature of $C$ is given by
\begin{equation}
\label{eq:kappa}
\kappa = -\dfrac{f_x^2 f_{yy} - 2 f_x f_y f_{xy} + f_y^2 f_{xx}}{|\nabla f|^3} = \dfrac{\begin{vmatrix}
\mathcal{H}(f) & (\nabla f)^T \\ 
\nabla f & 0
\end{vmatrix}}
{|\nabla f|^3},
\end{equation}
where $\mathcal{H}(f)$ is the Hessian matrix of $f$ (see \cite{Goldman2005} for example).

Let $g:\R^2\to\R$ be a smooth function, and set
\begin{equation}
\label{eq:Lagrange}
F:\R^2\times\R\to\R,\quad F(\x,\lambda)=g(\x)+\lambda f(\x).
\end{equation}
As mentioned above, if $\nabla f(\x)\ne 0$ and $\nabla g(\x) \ne 0$ then $\nabla F(\x,\lambda) = 0$ if and only if $C$ is tangent to $g(\x)=d$ at $\x$ for some $d$. 
It follows from that, when we take $\langle \x, \v \rangle$ $(\v\in S^1)$ and $|\x-\p|^2$ $(\p\in\R^2)$ as $g$ in \eqref{eq:Lagrange}, the condition $\nabla F=0$ holds if and only if $C$ is tangent to a line normal to $\v$ and to a circle centered at $\p$, respectively.

Consider
\begin{align*}
&H:\R^2\times\R\times S^1 \to \R,\quad H(\x,\lambda,\v) = \langle \x,\v \rangle + \lambda f(\x),\\
&D:\R^2\times\R\times \R^2 \to \R,\quad D(\x,\lambda,\p)=|\x-\p|^2 + \lambda f(\x),
\end{align*}
and define
\begin{align*}
h_{\v}(\x,\lambda)=H(\x,\lambda,\v),\quad
d_{\p}(\x,\lambda)=D(\x,\lambda,\p).
\end{align*}
According to the above reasons, we expect that $h_{\v}$ and $d_{\p}$ play the roles of the height and distance-squared functions on $C$, respectively. 

%%%%%%%%%%%%%%%%%%%%%%%%%%%%%%%%%%%%%%%
\subsection{Families of height functions}

\begin{theorem}
\label{thm:height}
Suppose that $f(\x_0)=0$ and $\nabla f(\x_0)\ne(0,0)$. Then
\begin{enumerate}
\item 
$\nabla h_{\v}=0$ at $(\x_0,\lambda_0)$ if and only if $\lambda_0=\pm 1/|\nabla f(\x_0)|$ and $\v=\mp \n(\x_0)$,  
\item 
$\nabla h_{\v}=0$ and $\det(\mathcal{H}(h_{\v}))=0$ at $(\x_0,\lambda_0)$ if and only if $\lambda_0=\pm 1/|\nabla f(\x_0)|$, $\v=\mp \n(\x_0)$ and $C$ has an inflection at $\x_0$.
 \end{enumerate}
\end{theorem}
\begin{proof}
We have $\nabla h_\v = (\v + \lambda \nabla f, f)$.
The assertion (1) is immediately deduced from this.
 
The Hessian matrix of $h_\v$ is given by
\[
\mathcal{H}(h_\v)
= 
\begin{pmatrix}
\lambda \mathcal{H}(f) & (\nabla f)^\mathrm{T} \\
\nabla f & 0
\end{pmatrix}.
\]
It follows from \eqref{eq:kappa} that
$$
\det(\mathcal{H}(h_\v)) = \lambda \kappa|\nabla f|^3 
$$
which establishes the assertion (2).
\end{proof}
Theorem \ref{thm:height} shows that $h_{\v}$ plays the role of the height function on $C$ along $\v$.

%%%%%%%%%%%%%%%%%%%%%%%%%%%%%%%%%%%%%%%%%%%%
\subsection{Families of distance-squared functions}


\begin{theorem}
\label{thm:distance}
Suppose that $f(\x_0)=0$ and $\nabla f(\x_0)\ne(0,0)$. Then
\begin{enumerate}
\item
$\nabla d_{\p}=0$ at $(\x_0,\lambda_0)$ if and only if $\p=\x_0+\lambda_0|\nabla f(\x_0)|\n(\x_0)/2$ holds.
\item
$\nabla d_{\p}=0$ and $\det(\mathcal{H}(d_{\p}))=0$ at $(\x_0,\lambda_0)$ if and only if $\lambda_0=2/(\kappa(\x_0)|\nabla f(\bm{x_0})|)$ and $\p=\x_0+\n(\x_0)/\kappa(\x_0)$. 
\item 
$\nabla d_{\p}=0$, $\det(\mathcal{H}(d_{\p}))=0$ and $\rank(J(\phi_d))<3$ at $(\x_0,\lambda_0)$ if and only if $\lambda_0=2/(\kappa(\x_0)|\nabla f(\bm{x_0})|)$, $\p=\x_0+\n(\x_0)/\kappa(\x_0)$ and $C$ has a vertex at $\x_0$, where
\[
\phi_d:\R^2\times\R\to\R^4,\quad \phi_d(\x,\lambda) =(\nabla d_{\p}(\x,\lambda),\det(\mathcal{H}(d_{\p}(\x,\lambda)))) 
\]
and $J(\phi_d)$ is the Jacobian matrix of $\phi_{d}$. 
\end{enumerate}
\end{theorem}

\begin{proof}
Since we have $\nabla d_\p=(2(\x-\p)+\lambda \nabla f,f)$, we obtain 
\begin{align*}
2(\x-\p)+\lambda |\nabla f|\,\n& =0, \\
\p&=\x+\dfrac{\lambda}2|\nabla f|\,\n,
\end{align*} 
which proves the assertion (1). 

The Hessian matrix of $d_\p$ is given by
$$
\mathcal{H}(d_\p) = 
\begin{pmatrix}
\begin{pmatrix}2&0\\0&2\end{pmatrix}+\lambda \mathcal{H}(f) & (\nabla f)^T\\
\nabla f & 0
\end{pmatrix},
$$
and thus we have
$$
\det(\mathcal{H}(d_\p))) = -2|\nabla f|^2+\lambda\left|\begin{matrix}\mathcal{H}(f)&(\nabla f)^T\\\nabla f&0\end{matrix}\right|.
$$
It follows from \eqref{eq:kappa} that 
$$
\det(\mathcal{H}(d_\p))=-2|\nabla f|^2+\lambda \kappa |\nabla f|^3.
$$
Therefore, 
\begin{align*}
-2|\nabla f|^2+\lambda \kappa |\nabla f|^3 & = 0 \\
\lambda & = \dfrac2{\kappa |\nabla f|},
\end{align*}
which proves the assertion (2). 

After a suitable rotation in the $xy$-plane if necessary, we may take
the $y$-axis as the normal direction to $C$ at $\x_0$. 
So we may assume that $f_x(\x_0)=0$ and $f_y(\x_0)\ne0$.  
Let $\lambda_0=2/(\kappa(\x_0)|\nabla f(\bm{x_0})|)$ and $\p=\x_0+\n(\x_0)/\kappa(\x_0)$ hold. 
Since $f_y(\x_0)\ne0$, by the implicit function theorem, the curve $f(x,y)=0$ can be expressed locally as the graph $y=y(x)$ satisfying $f(x,y(x))=0$ close to $\x_0$. 
It follows from \eqref{eq:kappa} that 
\begin{equation*}
\dfrac{d\kappa}{dx}(\x_0)=\dfrac{f_y(\x_0)^3(-f_{xxx}(\x_0)f_y(\x_0)+3f_{xx}(\x_0)f_{xy}(\x_0))}{|\nabla f(\x_0)|^5}.
\end{equation*}
Hence, $C$ has a vertex at $\x_0$ if and only if 
\begin{equation}
\label{eq:vertex}
-f_{xxx}(\x_0)f_y(\x_0)+3f_{xx}(\x_0)f_{xy}(\x_0)=0.
\end{equation}
We shall also calculate $J(\phi_d)$. 
A direct calculation of $J(\phi_d)$, similar to that for ridge points
of implicit surfaces in \cite{Hasegawa2017}, shows that $\rank(J(\phi_d))<3$ if and only if \eqref{eq:vertex} holds. 
\end{proof}


\begin{thebibliography}{99}
\bibitem[DT2018]{DT2018} F.~S.~Dias and F.~Tari, On vertices and inflections of plane curves, Journal of Singularities \textbf{17} (2018), 70--80.
\bibitem[Goldman2005]{Goldman2005} R.~Goldman, Curvature formulas for implicit curves and surfaces, Comput.\ Aided Geom.\ Design, \textbf{22} (2005), 632--658.
\bibitem[Hasegawa2017]{Hasegawa2017} M.~Hasegawa, Parabolic, ridge and sub-parabolic curves on implicit surfaces with singularities, Osaka J. Math. \textbf{54} (2017), 707--721.
\bibitem[IRRT2015]{IRRT2015} S.~Izumiya, M.~C.~Romero-Fuster, M.~A.~S.~Ruas and F.~Tari, Differential geometry from a singularity theory viewpoint, World Scientific Pub. Co Inc., Singapore (2015)
\bibitem[ST2017]{ST2017} M. Salarinoghabi and F. Tari, Flat and round singularity theory of plane curves, Q. J. Math., \textbf{68} (2017), 1289--1312.
\end{thebibliography}
\end{document}
